
En effet, soit $X$ dans $\mathbb P^N_k$. Quitte \`a remplacer $k$ par la cl\^oture alg\'ebrique d'une extension transcendante pure, ce qui est loisible gr\^ace \`a 1.6 (i), on peut supposer qu'il existe $r(N+1)$ \'el\'ements dans $k$, qui sont alg\'ebriquement ind\'ependants sur tous les $k_j$. Alors les $r$ sections hyperplanes ``g\'en\'eriques'' successives de $X$, d\'efinies au moyen de ces \'el\'ements, sont toujours des $k$-cycles positifs sp\'eciaux de degr\'e $d$.

Raisonnant par r\'ecurrence sur $r$, il reste donc \`a noter que lorsque $r=0$, alors $h^0({O}_X(-1))=d$, et lorsque $r\ge1$, alors $h^0({O}_X(-1))=0$. Pour ce dernier point, on peut supposer que $X$ lui-m\^eme provient d'un sch\'ema int\`egre $X'$, gr\^ace \`a l'injection ${O}_X\to \coprod{O}_{X_j}$ ; mais alors $h^0({O}_X(-1))=h^0({O}_{X'}(-1))$, ce qui est $0$ puisque $H^0({O}_{X'})$ est un corps.

\textbf{Corollaire 6.11.} (i) \emph{Soient $r,d$ deux entiers, $c_i=A_{r-i}(0,\ldots,0,d;0)$ pour $i=0,\ldots,r$ $(6.7\ \mathrm{(i)})$, et $p=P_r(c_0,\ldots,c_r)$ $(1.9)$. Soient $k$ un corps alg\'ebriquement clos, et $X$ un $k$-cycle positif sp\'ecial de dimension $\le r$ et degr\'e $\le d$, \`a polyn\^ome de Hilbert $\chi({O}_X(n))=\sum_{i=0}^{r}a_i\binom{n+i}{i}$. Alors pour $i=0,\ldots,r$, on a}
\[
   |a_i|\le c_i.
\]
\emph{De plus, $X$ s'immerge comme sous-sch\'ema de $\mathbb P^N_k$ avec $N=d(r+1)-1$, d\'efini par (au plus) $\binom{N+p}{N}$ \'equations de degr\'e $p$.}

(ii) (Chow) \emph{Soit $S$ un sch\'ema noeth\'erien. Lorsque $K$ parcourt les extensions alg\'ebriquement closes des $k(s)$, $s\in S$, les $K$-cycles positifs sp\'eciaux de dimension et degr\'e born\'es forment une ``famille limit\'ee'' (soit avec les $X$ consid\'er\'es comme sous-sch\'emas d'un $\mathbb P^N_S$ fix\'e, soit dans un sens abstrait \'evident).}

